% ===== 第六段：现有研究存在的问题（Gap分析）=====
\subsection{现有研究存在的问题}

综合上述四个方向的研究现状，现有工作在以下五个方面存在系统性不足，
这些不足构成本课题的直接研究动机：

\begin{enumerate}
  \item \textbf{G1：工具调用链风险度量体系缺失，组合性风险无法被现有per-call评估框架捕获。}
  现有基准（如AgentDojo\cite{balunovic2024agentdojo}）和防御方法以攻击成功率为核心指标，
  而对工具敏感度、权限扩散度、调用深度、不可逆操作比例、敏感数据暴露等过程性风险指标缺乏统一刻画。
  Chu等\cite{chu2026systematic}明确指出，语义权限提升发生时每次工具调用均经过授权，
  但授权动作的语义组合超越了预期授权范围——
  这类"序列合法但组合有害"的风险对现有以单次调用为粒度的访问控制系统完全不可见，
  也无法被任务完成率或攻击成功率等结果导向指标所量化。
  NIST AI RMF\cite{nist2023airmf}和MAESTRO\cite{huang2025maestro}均呼吁建立可测量的过程级风险指标，
  但具体的工具调用链度量方法尚未见于学术文献。

  \item \textbf{G2：Agent执行日志结构不统一，过程级审计缺乏可靠数据基础。}
  不同Agent框架记录的日志字段差异显著：部分框架只记录prompt与final answer，
  部分记录tool call与observation，鲜少完整记录计划步骤、参数来源、权限状态、
  工具返回对象、人工确认标记和跨Agent委托关系。
  Chu等\cite{chu2026systematic}指出，CoT轨迹被越来越多地用作审计日志，
  但其不可信性使基于CoT的审计存在根本缺陷；
  Dehghantanha等\cite{dehghantanha2026sokattacksurfaceagentic}则进一步指出，
  防篡改的执行追踪需要从设计层面将审计日志与Agent控制流分离。
  日志结构不统一导致后续安全度量和归因缺少标准化的数据基础。

  \item \textbf{G3：工具滥用检测局限于单步权限检查或最终结果判断，跨步骤组合风险路径识别不足。}
  现有防御机制（如权限白名单、提示注入过滤）主要在工具调用前进行单步权限检查，
  或在最终输出阶段判断违规内容，
  但对跨步骤、跨工具的风险路径识别不足——
  例如"读取敏感文件—摘要压缩—调用外部发送工具"或"查询数据库—写入共享记忆—下游Agent复用"
  等链式风险路径。
  de Witt等\cite{dewitt2026openchallenges}将这类攻击定义为链式攻击（chaining attacks），
  指出其利用的恰是"每次调用合法性与序列级意图之间的鸿沟"；
  Chu等\cite{chu2026systematic}指出，
  所有已发布的Agent安全基准均止步于T1/T2层面，
  跨步骤组合风险（T3）在结构上无法被现有评估框架测量。
  Zhang等\cite{zhang2024injectagent}和Liu等\cite{liu2024formalizing}的基准工作
  虽揭示了工具调用链的广泛脆弱性，但均以攻击成功率为最终指标，
  未提供面向组合风险路径的结构化检测方法。

  \item \textbf{G4：风险发现后的责任归因机制不足，审计结论缺乏可解释证据链。}
  当Agent系统出现不当工具调用或敏感数据外传时，
  仅输出"异常"二元标签无法满足实际审计需求。
  de Witt等\cite{dewitt2026openchallenges}指出，
  在代理链场景中，当多个Agent以管道形式协作时，
  确定哪个Agent执行了特定操作在技术上极为困难，
  导致现有审计日志不支持非否认性保证；
  Chu等\cite{chu2026systematic}将L1--L6层规避归因的攻击统一定义为L7治理层失效，
  说明缺乏归因机制是Agent系统安全体系中贯穿各层的系统性漏洞。
  实际应用需要回答：风险路径由哪些节点构成，关键触发工具是哪一个，
  参数来源于何处，哪个步骤对最终风险贡献最大——
  而SentinelAgent\cite{he2025sentinelagent}、AuthGraph\cite{wang2026aligning}等
  现有工具调用研究尚未提供满足上述需求的可解释归因方法。

  \item \textbf{G5：实验成本与工程复杂度制约研究落地，现有高成本评估框架难以低成本复现。}
  Deng等\cite{deng2025agentsunderthreat}指出，
  严格工具审计方法（如PrivacyAsst）相较标准Agent产生了高达100倍的额外计算开销，
  且仍未完全防止敏感信息泄露，说明高成本并不等价于高安全性；
  而Dziemian等\cite{dziemian2026vulnerable}的大规模红队竞赛虽覆盖了22个前沿Agent，
  但其框架依赖大量在线API调用，难以低成本复现。
  本课题需在较短周期内完成两个研究点、一个原型系统和会议论文投稿，
  必须设计实验成本可控、工具环境可模拟、标签可自动生成的低成本研究方案。
\end{enumerate}
